\section{Discusión}
\label{sec:discusion}

El diseño mantiene la separación de responsabilidades de la Entrega 1 y la
completa: la cola por niveles decide el orden de atención, el índice hash
responde dónde está cada solicitud, la lista doble permite retirarla sin
recorrer, y la pila registra el historial y permite revertirlo. Ninguna
estructura resuelve sola todas las operaciones; es su combinación, mantenida
consistente por la fachada, la que lleva todas las operaciones por
solicitud a tiempo constante (Tabla~\ref{tab:complejidad}).

La principal limitación de la Entrega 1, el costo lineal de la búsqueda, se
eliminó, pero a cambio de dos supuestos que conviene hacer explícitos. El
primero es que el costo $O(1)$ del índice es \emph{esperado}: depende de que
la función hash reparta bien los identificadores. Con los
\texttt{hashCode} de \texttt{String} y la mezcla de bits empleada, la cadena
más larga observada fue de 5 entradas con $10^4$ claves y de 6 con
$10^6$, pero un
conjunto de identificadores elegido para colisionar degradaría las
operaciones a $O(n)$. El segundo es el costo en memoria: cada pendiente
ocupa ahora un nodo con dos referencias y una entrada en el índice, unos
97 bytes en total (Sección~\ref{sec:resultados}).

El costo constante de la cola sigue dependiendo de que la urgencia tome
solo cinco valores. Si la prioridad dependiera de una magnitud continua,
como el tiempo de espera acumulado, habría que usar otra representación,
por ejemplo un \emph{heap} indexado con localizadores.

La operación de deshacer tiene una restricción deliberada: solo revierte
atenciones, no registros ni cancelaciones, y se rechaza si el identificador
fue reutilizado. Generalizarla a todas las operaciones exigiría apilar un
registro por cada tipo de acción, como en el patrón
\emph{Command}~\cite{gof}.

Finalmente, los datos existen solo durante la ejecución, lo cual es
consistente con el alcance definido.
